\documentclass[10pt,a4paper]{amsart}
\usepackage[margin=19mm]{geometry}
\usepackage{amsmath,amssymb,mathtools}
\usepackage[scr=rsfs,cal=euler]{mathalfa}
\usepackage{xfrac}
\usepackage[unicode,hidelinks,bookmarks=false]{hyperref}
\newcommand{\ie}{i.e.\ }
\newcommand{\cf}{cf.\ }
\newcommand{\resp}{respectively\ }
\newcommand{\Subsection}{\subsection}
\providecommand{\nobreakdash}{\nobreak\hbox{-}}
\setlength{\emergencystretch}{2em}
\multlinegap0pt
\newcommand{\CA}[1]{$(\textsc{c}_{#1})$}
\usepackage{amssymb}
\usepackage{xcolor}
\usepackage{algorithm}\usepackage{algpseudocode}
\usepackage{tikz}
\newtheorem{lemma}{Lemma}
\newtheorem{claim}{Claim}
\newtheorem{claimproofnew}{Claimproofnew}
\newcommand{\E}{{\embed{\WW}}}
\newcommand{\G}{\mathcal{G}}
\newcommand{\M}{\mathcal{M}}
\newcommand{\Pres}[1]{\mathcal{P}_{#1}}
\newcommand{\R}{\mathbb{K}}
\newcommand{\Sof}[1]{\mathcal{S}(#1)}
\newcommand{\Ssub}[1]{\mathcal{S}_{#1}}
\newcommand{\U}{U}
\newcommand{\UU}{\mathcal{C}}
\newcommand{\V}{Y}
\newcommand{\VV}{\mathcal{B}}
\newcommand{\Vs}{B}
\newcommand{\W}{A}
\newcommand{\WW}{\mathcal{A}}
\newcommand{\age}[1]{\text{\sc Age}(#1)}
\newcommand{\defeq}{\stackrel{\text{\tiny def}}{=}}
\newcommand{\embed}[1]{\text{\sc Emb}_{#1}}
\newcommand{\embset}[1]{\restr \E #1}
\newcommand{\embsetprod}[2]{#1\boxtimes #2}
\newcommand{\minupperb}[2]{#1 \uparrow #2}
\newcommand{\polyring}[2]{#1[#2]}
\newcommand{\renorder}{\leadsto}
\newcommand{\restr}[2]{{#1}{|#2}}
\newcommand{\set}[1]{\left\{#1\right\}}
\newcommand{\setof}[2]{\left\{#1 \, \middle\vert \, #2\right\}}
\newcommand{\subseteqfin}{\subseteq_\text{fin}}
\newcommand{\vars}[1]{\text{\sc vars}(#1)}
\newcommand{\wpo}{\sqsubseteq}
\title{Equivariant ideals of polynomials}
\author{Arka Ghosh, S{\l}awomir Lasota}
\date{Source: arXiv:2402.17604v3 (CC BY 4.0)}
\begin{document}
\maketitle

\noindent\emph{Source and licence.} Equivariant ideals of polynomials. Version arXiv:2402.17604v3 (CC BY 4.0), distributed by arXiv under the Creative Commons Attribution 4.0 International licence, http://creativecommons.org/licenses/by/4.0/, as recorded in the arXiv licence metadata for this exact version and shown on the abstract page https://arxiv.org/abs/2402.17604v3. Full licence notice: \texttt{licenses/arxiv-2402.17604-CC-BY-4.0.txt}.\par\smallskip

\noindent\textit{Editorial scope.} Counted unit: the article's lemma finite (source0.tex lines 3006--3013). The article's complete original proof of it is delivered with it (source0.tex lines 3015--3171, one \texttt{proof} environment of 157 lines). No mathematics is written, repaired or re-derived: the statement and the proof are the article's own lines, in the article's own notation, and no retained line is substituted or re-flowed.  Retained uncounted as the source-local context the proof needs: lines 2137--2139; lines 2156--2159; lines 2989--2991.  Nothing was omitted from any retained span, so the package carries no omission record.  No retained line contains a citation, so the wrapper carries no bibliography and no bibliography entry.  No retained line contains a figure environment, an image-inclusion command or drawing code, so no image is delivered and the package declares no asset dependency.  Editorial formatting only, in addition: an AMS \texttt{amsart} preamble that restates the article's own declarations and macros used by the retained lines, namely the environments \texttt{lemma}, \texttt{claim}, \texttt{claimproofnew} on separate counters, together with the standard packages those lines need.  The retained environments print in this document as Lemma 1, Claim 1, Claimproofnew 1 in order of appearance, whereas the article prints the same items with the numbering of its own full document.  The complete native source of the article as distributed in the arXiv e-print for this exact version is bundled unchanged as \texttt{sources/155220/source0.tex}.
\begin{lemma} \label{lem:restr}
$\restr \iota {\vars f} = \restr {\iota'} {\vars f} \implies \iota(f) = \iota'(f)$.
\end{lemma}

\begin{lemma} \label{lem:embsetprod}
For $\Vs, \Vs'\in\age \WW$,  $\embsetprod \Vs  {\Vs'}$ is orbit-finite, namely
$\embsetprod \Vs  {\Vs'} = \E(F) \defeq \setof{(\kappa\circ\pi,\kappa\circ\pi')}{\kappa\in\E, (\pi, \pi')\in F}$ for some finite  $F\subseteq \embsetprod \Vs {\Vs'}$.
\end{lemma}

%\begin{align*}%\label{eq:commutes}
$\Sof {\kappa(f), \kappa(g)} \ = \ \kappa(\Sof {f,g})$
%\end{align*}

\begin{lemma}[Computable presentation] \label{lem:min finite}
For every finite set of polynomials $\G \subseteq \polyring{\R}{\WW}$, the set $\Ssub\G$ 
is orbit-finite, namely
\[
\Ssub\G \ = \ \E(\Pres\G) \ \defeq \ \setof{\iota(f)}{\iota\in\E, f\in\Pres\G}
\]
for some presentation $\Pres\G \subseteqfin \Ssub\G$, which is moreover computable.
\end{lemma}

\begin{proof}
We start by showing orbit-finiteness of $\Ssub\G$.
As $\G$ is finite and $\Ssub \G = \bigcup_{g,g'\in\G} \Ssub {g,g'}$, where
\[
\Ssub {g,g'} \ = \  \setof{\Sof{\iota(g),\iota'(g')}}{\iota, \iota' \in \E, \ \iota(g)\neq\iota'(g')},
\]
it is enough to prove that $\Ssub {g,g'}$ is orbit-finite for every two fixed polynomials $g,g'\in\G$.
Let $\Vs = \vars g$ and $\Vs' = \vars {g'}$, and
let $F_{g,g'}$ be a finite presentation of $\embsetprod {\Vs} {\Vs'}$
(recall Lemma \ref{lem:embsetprod}).
%Let $E$ denote the set of all pairs of local embeddings $(\iota, \iota')$ that map $\VV$ and $\VV'$ to arbitrary but
%the same structure $\UU$ in F:
%\[
%E \ = \ \setof{(\iota, \iota')}{\UU\in F, \ \iota\in\embset\VV{\UU}, \ \iota'\in\embset{\VV'}{\UU}}.
%\]
Relying on Lemma \ref{lem:restr},
for a local embedding $\pi\in\embset\Vs$ we write $\pi(g)$ to mean the renaming of $g$
by \emph{any} extension of $\pi$ to an embedding of $\WW$.
Likewise we write $\pi'(g')$, for $\pi'\in\embset{\Vs'}$:
\[
\Ssub {g,g'} \ = \  \setof{\Sof{\pi(g),\pi'(g')}}{(\pi, \pi') \in \embsetprod \Vs {\Vs'}, \ \pi(g)\neq\pi'(g')}.
\]
Finally, we define $\Pres{g,g'}\subseteq \Ssub {g,g'}$ as the subset of $\Ssub {g,g'}$ obtained by restricting 
to local pairs of embeddings $(\pi, \pi')$ belonging to $F_{g,g'}$:
\[
\Pres{g,g'} \ = \ \setof{\Sof {\pi(g), \pi'(g')}}{(\pi,\pi') \in F_{g,g'}, \ \pi(g)\neq\pi'(g')}.
\]
As $F$ is finite, $\Pres{g,g'}$ is finite as well, 
and orbit-finiteness of $\Ssub {g,g'}$ follows once we prove that
$\Pres{g,g'}$ is a presentation of $\Ssub{g,g'}$:
%
\begin{claim}
$\Ssub {g,g'} = \E(\Pres{g,g'})$.
% \defeq \setof{\kappa(f)}{\kappa\in\E, \, f\in \Pres{\set{g,g'}}}$.
\end{claim}
\begin{claimproofnew}
For proving the inclusion
$\E(\Pres{g,g'}) \subseteq \Ssub {g,g'}$,
we take any polynomial $\kappa(f)\in \E(\Pres{g,g'})$, where $\kappa \in \E$ and $f = \Sof {\pi(g), \pi'(g')} \in \Pres{g,g'}$,
and argue that $\kappa(f) \in \Ssub {g,g'}$.
As action of embeddings commutes with S-polynomials, 
we have
\[
\kappa(\Sof {\pi(g), \pi'(g')}) \ = \ \Sof {\kappa(\pi(g)), \kappa(\pi'(g'))}.
\]
As $\kappa$ is injective, 
knowing $\pi(g)\neq\pi'(g')$, we also know $\kappa(\pi(g))\neq\kappa(\pi'(g'))$,
and therefore $\kappa(f) \in \Ssub {g,g'}$.

The opposite inclusion $\Ssub {g,g'} \subseteq \E(\Pres{g,g'})$ is shown similarly.
We take any polynomial $\Sof {\pi(g), \pi'(g')} \in \Ssub {g,g'}$.
As $F_{g,g'}$ is a presentation of $\embsetprod \Vs{\Vs'}$, there is some 
$(\overline\pi, \overline\pi')\in F_{g,g'}$ and
$\kappa\in\E$ such that
$(\pi,\pi') = (\kappa\circ\overline\pi, \kappa\circ\overline\pi')$.
Knowing $\kappa(\pi(g))\neq\kappa(\pi'(g'))$, we also know $\pi(g)\neq\pi'(g')$,
and therefore $\Sof {\overline\pi(g), \overline\pi'(g')} \in \Pres{g,g'}$.
As action of embeddings commutes with S-polynomials, we get
\[
\Sof {\pi(g), \pi'(g')} \ = \ 
\kappa(\Sof {\overline\pi(g), \overline\pi'(g')}),
\]
and in consequence $\Sof {\pi(g), \pi'(g')}\in\E(\Pres{g,g'})$.
\end{claimproofnew}

%\smallskip

Union of all the sets $\Pres{g,g'}$ yields therefore a presentation of $\Ssub \G$:
\[
\Pres{\G} \ \defeq \ \bigcup_{g,g'\in\G}\Pres{g,g'}.
\]
It is moreover computable:
for every $g,g'\in\G$ the finite set $F_{g,g'}$ is computable, by assumption \CA 2,
and therefore the set $\Pres{g,g'}$ is computable as well.
%
%\bigskip
%
%Let $g,g'\in\G$ be arbitrary two fixed polynomials from $\G$, and
%let $\V,\V'\subseteqfin \W$ denote the sets of variables appearing in $g$ and $g'$, respectively.
%As $\G$ is finite, it is enough to argue that the following set is finite and computable:
%\begin{align}\label{eq:Mdef}
%\M \ = \ \min_{\renorder} \setof{\Sof{\iota(g),\iota'(g')}}{\iota, \iota' \in \E, \ \iota(g)\neq\iota'(g')}.
%\end{align}
%%We say that a pair of embeddings $\iota,\iota'\in\E$ \emph{factorizes} through $\kappa\in\E$
%%if 
%%\[
%%\iota = \kappa \circ \bar \iota \qquad
%%\iota' = \kappa \circ \bar \iota'
%%\]
%%for some $\bar\iota, \bar\iota'\in\E$.
%We recall that action of embeddings $\kappa\in\E$ commutes with S-polynomials, 
%as in equality \eqref{eq:commutes}
%which we can instantiate to get:
%\begin{align}\label{eq:comm}
%\Sof {\kappa(\iota(g)), \kappa(\iota'(g'))} \ = \ 
%\kappa(\Sof {\iota(g),\iota'(g')}).
%\end{align}
%Furthermore, whenever a pair $(\restr\iota\V, \restr{\iota'}{\V'})$ factorizes through another such pair
%$(\restr{\bar\iota}\V, \restr{\bar\iota'}{\V'})$, namely
%\[
%\restr{\iota} \V = \restr{(\kappa \circ \bar \iota)} \V \qquad
%\restr{\iota'} {\V'} = \restr{(\kappa \circ \bar \iota')} {\V'},
%\]
%knowing that $\iota(g)$ is determined uniquely by $\restr\iota \V$, and likewise
%$\iota'(g')$ is determined uniquely by $\restr{\iota'}{\V'}$,
%we may rewrite \eqref{eq:comm} to get:
%\[
%\Sof {\iota(g), \iota'(g')}
% \ = \ 
% \kappa(\Sof {\bar\iota(g),\bar\iota'(g')}).
%\]
%Therefore,
%whenever a pair $(\restr\iota\V, \restr{\iota'}{\V'})$ factorizes through another such pair
%$(\restr{\bar\iota}\V, \restr{\bar\iota'}{\V'})$, the S-polynomial
%$\Sof {\iota(g), \iota'(g')}$ is not in $\M$.
%Thus
%$\M$ contains exactly those S-polynomials $\Sof{\iota(g),\iota'(g')}$ where the pair
%$(\restr\iota\V, \restr{\iota'}{\V'})$ factorizes through no other pair, i.e., is 
%where the pair $(\iota,\iota')$ is $(\V,\V')$-minimal:
%\[
%\M \ = \ \setof{\Sof{\iota(g),\iota'(g')}}{(\iota, \iota') \in (\E)^2 \text{ a $(\V,\V')$-minimal pair}, \, \iota(g)\neq\iota'(g')}.
%\]
%It therefore suffices to show finiteness and computability of the following set of pairs of polynomials:
%\[
%%\M^{(1)} \ = \ 
%\setof{(\iota(g),\iota'(g'))}{(\iota, \iota') \in (\E)^2 \text{ a $(\V,\V')$-minimal pair}, \, \iota(g)\neq\iota'(g')},
%\]
%which is implied by finiteness and computability of its superset where the last restriction is dropped:
%\[
%%\M^{(2)} \ = \ 
%\setof{(\iota(g),\iota'(g'))}{(\iota, \iota') \in (\E)^2 \text{ a $(\V,\V')$-minimal pair}}.
%\]
%As $\iota(g)$ is determined uniquely by $\restr\iota \V$, and likewise
%$\iota'(g')$ is determined uniquely by $\restr{\iota'}{\V'}$,
%it suffices to show finiteness and computability of the following set of
%pairs of sets of variables:
%\[
%%\M^{(3)} \ = \ 
%\setof{(\restr\iota\V,\restr{\iota'}{\V'})}{(\iota, \iota') \in (\E)^2 \text{ a $(\V,\V')$-minimal pair}}.
%\]
%This set is exactly $\minupperb \V {\V'}$ and hence is finite and computable 
%by assumption \CA 2.
%%
%%Furthermore, as $\V$ or $\V'$ embeds into $\U=\iota(\V) \cup \iota'(\V')$
%%in finitely many ways (once $\U$ is fixed), 
%%and one can compute all these embeddings by the computability assumption \CA 1,
%%it is enough to show finiteness and computability of the following set of sets of variables:
%%\[
%%\M^{(4)} \ = \ \setof{(\iota(\V) \cup \iota'(\V'))}{(\iota, \iota') \in \E^2 \text{ a minimal pair}}.
%%\]
%%To complete the proof we observe that $M^{(4)}$ is exactly the set of minimal upper bounds of 
%%$\V$ and $\V'$ with respect to the well partial order $\wpo$,
%%namely $\M^{(4)} = \minupperb \V {\V'}$
%%(and therefore $\M^{(4)}$ is finite), 
%%and use %the fact that the latter set is finite, and 
%%the computability assumption \CA 2 that $\minupperb \V {\V'}$ is computable.
\end{proof}

\end{document}
